\section{Tangential curves and the exact amount of curve data}
\label{sec:curves}

Straight rays are only the simplest restrictions of the ordered germ.
A curve can be tangent to one or several prefix hyperplanes, producing
a higher pole order. The canonical decomposition gives a finite
calculus for all such cases, including a sharp statement about which
Taylor coefficients of the curve can affect the answer.

\subsection{A coefficient formula for every admissible analytic curve}

Let $u_i(\eps)$ be analytic near zero, with $u_i(0)=0$, and let
\begin{equation}\label{curve:prefix}
 P_j(\eps)=\sum_{i=1}^j u_i(\eps).
\end{equation}
Assume that no $P_j$ is identically zero. There are unique integers
$\nu_j\ge1$ and analytic units $p_j$ such that
\begin{equation}\label{curve:units}
 P_j(\eps)=\eps^{\nu_j}p_j(\eps),\qquad p_j(0)\ne0.
\end{equation}
Set $M_0=0$ and $M_k=\nu_1+\cdots+\nu_k$.

\begin{theorem}[Exact curve finite part and sharp jet order]
\label{curve:main}
The composed ordered function has a pole of exactly order $M_d$, with
leading coefficient $\prod_{j=1}^d p_j(0)^{-1}$. Its constant term is
\begin{equation}\label{curve:FP}
 \boxed{\displaystyle
 \FP_{\eps=0}Z_d(1+u_1(\eps),\ldots,1+u_d(\eps);a)
 =\sum_{k=0}^d[\eps^{M_k}]
 \frac{R_{d-k}(u_{k+1}(\eps),\ldots,u_d(\eps);a)}
      {p_1(\eps)\cdots p_k(\eps)}.}
\end{equation}
Regular-germ Taylor coefficients through total degree $M_d$ and curve
Taylor coefficients through order
\begin{equation}\label{curve:order}
 J=M_d+\max_{1\le j\le d}\nu_j
\end{equation}
suffice. For the class of curves with the specified prefix orders,
the bound $J$ cannot in general be decreased.
\end{theorem}

\begin{proof}
Insert \eqref{curve:units} into \eqref{ord:canonicalformula}. Its $k$th
summand is $\eps^{-M_k}$ times a holomorphic function. The constant is
therefore precisely the coefficient on the right of \eqref{curve:FP}.
Because $M_0<M_1<\cdots<M_d$, only the $k=d$ term can have pole order
$M_d$. Its numerator is $R_0=1$, so its leading coefficient is nonzero
and equals the asserted product.

In the $k$th numerator, each $u_i$ vanishes at least to first order.
Coefficients through order $M_k$ thus use only regular-germ monomials
of total degree at most $M_k$, and only curve coefficients through
that order. To invert $p_j=P_j/\eps^{\nu_j}$ through degree $M_k$
requires $P_j$ through degree $\nu_j+M_k$. These degrees are all at
most \eqref{curve:order}, proving sufficiency.

For necessity, fix an index $j$ and perturb only that prefix by
\begin{equation}\label{curve:perturb}
 P_j(\eps)\longmapsto P_j(\eps)+
                       \delta\eps^{\nu_j+M_d}.
\end{equation}
This is realized by opposite changes to $u_j,u_{j+1}$ when $j<d$,
or by a change to $u_d$ when $j=d$. All prefix orders and all other
prefixes stay fixed. The affected unit first changes in degree $M_d$.
Every denominator change in a term with $k<d$ starts after its constant
degree $M_k$; numerator changes start even later. In the $k=d$ term,
ordinary inversion gives the exact change
\begin{equation}\label{curve:sharpchange}
 \Delta\FP=-\frac{\delta}
 {p_j(0)\prod_{\ell=1}^d p_\ell(0)}.
\end{equation}
Choosing $j$ with maximal $\nu_j$ changes no curve coefficient below
order $J$ but changes the finite part by a nonzero multiple of
$\delta$. This proves sharpness.
\end{proof}

The theorem does not assign a restriction to a curve lying identically
in a polar hyperplane. Such a restriction needs an additional
operation, for example first removing a specified polar germ. A
finite algebraic value obtained by continuing a slope formula is not
by itself that extra prescription.

\subsection{Transverse curves and a depth-three fourth-jet formula}

When every prefix is transverse, $\nu_j=1$, Theorem~\ref{curve:main}
requires the curve jet through degree $d+1$. Thus a quadratic curve
description is insufficient even at depth two, and a cubic description
is insufficient in general at depth three.

An explicit example is a common reparametrization of a fixed ray.
Suppose
\[
 Z_3(1+c_1t,1+c_2t,1+c_3t;a)
 =A_{-3}t^{-3}+A_{-2}t^{-2}+A_{-1}t^{-1}+A_0+O(t),
\]
where all coefficients are given by
\eqref{ord:threeFP}--\eqref{ord:threepoles}. Put
\[
 t=h(\eps)=\eps+\alpha\eps^2+\beta\eps^3+
                         \chi\eps^4+O(\eps^5).
\]
Then exact series inversion yields
\begin{equation}\label{curve:cubicreparam}
 \boxed{\begin{aligned}
 &\FP_{\eps=0}Z_3(1+c_1h(\eps),1+c_2h(\eps),1+c_3h(\eps);a)\\
 &\quad=A_0-\alpha A_{-1}+(3\alpha^2-2\beta)A_{-2}\\
 &\qquad\quad+(-10\alpha^3+12\alpha\beta-3\chi)A_{-3}.
 \end{aligned}}
\end{equation}
For a transverse ray, $A_{-3}\ne0$, so the fourth coefficient $\chi$
has an unavoidable effect. The identity follows by taking the
constant coefficients of $h^{-1},h^{-2},h^{-3}$; their respective
values are $-\alpha$, $3\alpha^2-2\beta$, and
$-10\alpha^3+12\alpha\beta-3\chi$.

At the equal-slope ray $c_1=c_2=c_3=1$, one may insert
\[
 A_{-3}=\frac16,\qquad A_{-2}=\frac{\gamma_0(a)}2,\qquad
 A_{-1}=\frac{\gamma_0(a)^2-\gamma_1(a)-\zeta(2,a)}2,
\]
and $A_0=\mathcal F_3(1,1;a)$. This supplies a completely explicit
curved depth-three identity using ordinary Hurwitz data.

\subsection{A tangency example of arbitrarily high order}

At depth two take
\[
 u_1(\eps)=\eps,\qquad u_2(\eps)=-\eps+\eps^r,
 \qquad r\ge2.
\]
Here $(\nu_1,\nu_2)=(1,r)$, so the pole order is $r+1$ and the
sharp jet order is $2r+1$. The canonical formula simplifies exactly to
\[
 Z_2=\eps^{-r-1}+\frac{R_1(-\eps+\eps^r;a)}\eps
                  +R_2(\eps,-\eps+\eps^r;a).
\]
Consequently
\begin{equation}\label{curve:tangentexample}
 \FP Z_2=B_2(a)+\gamma_1(a).
\end{equation}
Now change $u_2$ by $\delta\eps^{2r+1}$. This leaves all lower curve
coefficients unchanged, but the first displayed term contributes an
additional constant $-\delta$. Thus the very short unperturbed answer
does not reduce the amount of curve information needed in general.

\subsection{Dependence on the ratio of straight slopes}

For $q\ne0$, put $\lambda=p/q$ and write
$\mathcal F_d(\lambda;a)=\mathcal F_d(p,q;a)$. Theorem
\ref{ord:oneslope} gives the rational function
\begin{equation}\label{curve:ratio}
 \mathcal F_d(\lambda;a)=B_d(a)+
  \sum_{k=1}^{d-1}\frac{[y^{d-k}t^k]F_a(y,t)}{(\lambda)_k}.
\end{equation}
It has at most simple poles at $0,-1,\ldots,2-d$. In particular,
for $d\ge2$,
\begin{equation}\label{curve:stieltjesresidue}
 \boxed{\displaystyle
 \operatorname*{Res}_{\lambda=2-d}\mathcal F_d(\lambda;a)
 =-\frac{\gamma_{d-1}(a)}{(d-1)!(d-2)!}.}
\end{equation}
Only the $k=d-1$ summand can contribute to this residue; its numerator
is $(-1)^{d-1}\gamma_{d-1}(a)/(d-1)!$, and the remaining factors in
the denominator have product $(-1)^{d-2}(d-2)!$.
Equation~\eqref{curve:stieltjesresidue} is an identity of the rational
slope continuation. It is not a value of the original multiple zeta
function on the excluded polar ray. It gives a useful inverse
relation: a higher Stieltjes constant can be recovered from the
slope dependence of the ordered finite-part family.

In fact the entire coefficient family has an exact inverse in terms
of these residues. For $d\ge2$, put
\[
 A_k=[y^{d-k}t^k]F_a(y,t),\qquad
 \rho_j=\operatorname*{Res}_{\lambda=-j}\mathcal F_d(\lambda;a),
 \quad 0\le j\le d-2.
\]
Partial fractions in \eqref{curve:ratio} give
\begin{equation}\label{curve:residuetransform}
 \rho_j=\frac{(-1)^j}{j!}
          \sum_{k=j+1}^{d-1}\frac{A_k}{(k-1-j)!},\qquad
 \boxed{A_k=(-1)^{k-1}\sum_{j=k-1}^{d-2}
                   \frac{j!}{(j-k+1)!}\rho_j.}
\end{equation}
To verify the inverse, set $b_j=(-1)^j j!\rho_j$. The first transform
is the upper-triangular convolution with $1/r!$; convolution with
$(-1)^r/r!$ inverts it, because $e^z e^{-z}=1$. This proves the second
formula and recovers \eqref{curve:stieltjesresidue} at $k=d-1$.

\subsection{Recovering every regular jet from directional finite parts}

The one-distinguished-slope family determines diagonal combinations
of regular coefficients. Allowing all independent slopes recovers
each coefficient that can appear, with no loss of information.

Write $H_{r,m}(\boldsymbol v;a)=[t^m]R_r(t\boldsymbol v;a)$ for the
homogeneous degree-$m$ Taylor polynomial, and let
\[
 \mathcal G_d(\boldsymbol c;a)=\FP_{\eps=0}
 Z_d(1+c_1\eps,\ldots,1+c_d\eps;a).
\]
For $P_j=c_1+\cdots+c_j\ne0$, the canonical formula gives
\begin{equation}\label{curve:generalray}
 \mathcal G_d(\boldsymbol c;a)=B_d(a)+
 \sum_{k=1}^{d-1}
 \frac{H_{d-k,k}(c_{k+1},\ldots,c_d;a)}{P_1\cdots P_k}.
\end{equation}

\begin{corollary}[Exact reconstruction from directional data]
\label{curve:reconstruction}
The rational function $\mathcal G_d$ uniquely determines every
$H_{d-k,k}$, $1\le k\le d-1$, and its constant $B_d$. In particular,
the full collection of directional finite parts at all depths
determines every regular germ $R_r$: its homogeneous degree-$m$ part
is recovered at depth $r+m$.
\end{corollary}

\begin{proof}
Use the independent prefix coordinates $P_1,\ldots,P_d$ and descend
from $k=d-1$ to $k=1$. Suppose all higher-index summands in
\eqref{curve:generalray} have been recovered and subtracted. Take the
coefficient of $P_1^{-1}\cdots P_k^{-1}$ in the remainder, treating
$P_{k+1},\ldots,P_d$ as generic parameters. The lower-index summands
are holomorphic in at least one of these first $k$ variables and
contribute zero. The surviving coefficient is
\[
 H_{d-k,k}(P_{k+1},P_{k+2}-P_{k+1},\ldots,P_d-P_{d-1};a).
\]
This invertible change of the remaining variables recovers the
entire polynomial $H_{d-k,k}$. Subtracting all recovered terms leaves
$B_d$. For a general $R_r$ and $m\ge1$, choose $d=r+m$ and $k=m$.
The degree-zero part is $R_r(0)=B_r$, already known from the Gamma
quotient. No higher Taylor degree than those in
\eqref{curve:generalray} can affect a finite part at the fixed depth.
\end{proof}
